\section{GEMM vs. GEMV}

\begin{table}[H]
\centering
\caption{Prefill vs. Decode 计算特性对比}
\begin{tabular}{lcc}
\toprule
& \textbf{Prefill} & \textbf{Decode} \\
\midrule
核心操作 & GEMM（矩阵$\times$矩阵） & GEMV（向量$\times$矩阵） \\
Q 长度 & $N$（prompt 长度） & 1 \\
瓶颈类型 & 计算密集型 & 内存带宽密集型 \\
GPU 利用 & 算力充分利用 & 算力严重闲置 \\
优化方向 & 更强的 FLOPs & 更大的内存带宽 \\
延迟特征 & TTFT（首 token 延迟） & TPS（每秒 token 数） \\
\bottomrule
\end{tabular}
\end{table}

\begin{warningbox}{Decode 阶段的效率困境}
Decode 阶段每步只生成一个 token，但需要加载整个模型权重和 KV-Cache。GPU 的算力大量闲置，内存带宽成为瓶颈。这就是为什么 Apple M5 Max 的统一内存架构（高带宽）在 Decode 阶段表现突出。
\end{warningbox}

